\chapter{Design Patterns and Industry Standards}
\label{chap:ch4}

\par Modern software development uses well-established design patterns to keep complexity under control and make code easier to maintain over time. Although neither C nor Rust includes a built-in concept of classes, both languages provide effective mechanisms for implementing object-oriented design patterns. This chapter explores how each language applies five fundamental patterns that form the foundation of the Gentlix Bank application.

\section{Encapsulation}
\label{sec:ch4sec1}

Encapsulation hides implementation details from client code, exposing only a
controlled interface. This principle protects internal invariants and permits
implementation changes without breaking dependent code.

In C, encapsulation is achieved through the opaque pointer pattern. A header file
declares a forward struct declaration without revealing its members, and all
operations are exposed as functions accepting and returning pointers to this
incomplete type. For the transaction service, the header declares:

\begin{lstlisting}[language=C, caption={transaction\_service.h --- Opaque pointer
hides internal structure. Client code sees only the forward declaration and
the public function signatures; fields are completely invisible.}]
// Forward declaration only
typedef struct TransactionService TransactionService;

// Constructor / Destructor
TransactionService *transaction_service_new(
    TransactionRepository tx_repo,
    TxAccountRepository   account_repo);

void transaction_service_free(TransactionService *svc);

// Public methods
bool transaction_service_record_transaction(
    TransactionService             *svc,
    const RecordTransactionCommand *cmd,
    TransactionId                  *out_id,
    ServiceError                   *err);
\end{lstlisting}

The actual definition resides in the implementation file where client code cannot
access it:

\begin{lstlisting}[language=C, caption={transaction\_service.c --- The concrete
struct is defined only in the implementation file. Client code cannot read or
modify \texttt{tx\_repo} and \texttt{account\_repo} directly; all access goes
through the public functions declared in the header.}]
// Visible only inside this translation unit
struct TransactionService {
    TransactionRepository tx_repo;      // polymorphic vtable + ctx
    TxAccountRepository   account_repo; // polymorphic vtable + ctx
};

TransactionService *transaction_service_new(
    TransactionRepository tx_repo,
    TxAccountRepository   account_repo)
{
    TransactionService *svc = malloc(sizeof *svc);
    if (!svc) return NULL;
    svc->tx_repo      = tx_repo;
    svc->account_repo = account_repo;
    return svc;   // caller owns this pointer; must call _free()
}

void transaction_service_free(TransactionService *svc) { free(svc); }
\end{lstlisting}

This technique, coined by Schreiner~\cite{Schreiner1993} as \emph{information hiding
through incomplete types}, ensures that no client code can depend on internal layout.
Even if the struct gains new fields or changes internally, only the implementation 
file needs to be recompiled --- client code is not affected.

Rust's module system provides more granular control than C's header file
boundary. The \texttt{pub} keyword governs visibility at the field, function,
and module level independently --- a struct can be public while its fields
remain private. The compiler enforces these boundaries: code outside the
defining module cannot access or construct a type whose fields are private,
regardless of how the struct is named or where it is used.

In \texttt{TransactionService}, both repository handles are private fields.
The only way to build an instance is through the public \texttt{new}
constructor, which receives the repositories as arguments:

\begin{lstlisting}[language=Rust, caption={transaction.rs (service) --- Fields
are private by default. Outside this module, the only way to create a
\texttt{TransactionService} is through \texttt{new}.}]
pub struct TransactionService<T, A>
where
    T: TransactionRepository,
    A: AccountRepository,
{
    tx_repo: Arc<T>,      // not accessible outside this module
    account_repo: Arc<A>, // not accessible outside this module
}

impl<T, A> TransactionService<T, A>
where
    T: TransactionRepository,
    A: AccountRepository,
{
    pub fn new(tx_repo: Arc<T>, account_repo: Arc<A>) -> Self {
        Self { tx_repo, account_repo }
    }
}
\end{lstlisting}

As Palmieri notes~\cite{Palmieri2022}, making incorrect usage patterns
\emph{unrepresentable by construction} is the central idea of type-driven
development in Rust: ``we made it impossible for an instance of
\texttt{SubscriberName} to violate those constraints [\ldots] it will not
compile.'' The same principle applies here --- a caller who attempts to
bypass \texttt{new} and set \texttt{tx\_repo} directly receives a
compile-time error, not a runtime failure.

Unlike C's all-or-nothing approach --- where encapsulation depends on
keeping the struct definition out of the header file --- Rust permits
fine-grained control at the field level within a single module. Two fields
of the same struct can have different visibility, and the compiler verifies
every access point across the entire codebase automatically.

\section{Methods and impl Blocks}
\label{sec:ch4sec2}

A method is a function that belongs to a specific type. Without this, a codebase 
can end up as a long list of functions with no clear indication of 
which type they belong to.

In C, methods do not exist as a language feature. Instead, programmers
simulate them through a naming convention: every function that operates on
a TransactionService begins with the prefix
\texttt{transaction\_service\_}, and its first parameter is always a pointer
to the struct. This pointer plays the role of \texttt{self} in
object-oriented languages. Schreiner~\cite{Schreiner1993} describes this as
``messages to objects'' --- the first argument identifies the receiver, and
the function name identifies the operation.

The constructor and destructor follow the same convention:

\begin{lstlisting}[language=C, caption={transaction\_service.c ---
Constructor and destructor follow the \texttt{type\_verb} naming convention.
The \texttt{svc} parameter is the implicit ``self'' of the method.}]
TransactionService *transaction_service_new(
    TransactionRepository tx_repo,
    TxAccountRepository   account_repo)
{
    TransactionService *svc = malloc(sizeof *svc);
    if (!svc) return NULL;
    svc->tx_repo      = tx_repo;
    svc->account_repo = account_repo;
    return svc;   /* caller owns this; must call _free() */
}

void transaction_service_free(TransactionService *svc) { free(svc); }
\end{lstlisting}

A more complex method, such as recording a transaction, follows the same
pattern. The \texttt{svc} pointer gives access to the struct's fields via
\texttt{->}, and all error reporting goes through output parameters.
This approach is purely conventional --- the compiler does not know that
\texttt{transaction\_service\_record\_transaction} belongs to
\texttt{TransactionService}, and nothing prevents passing the wrong pointer
or omitting the prefix entirely:

\noindent\begin{minipage}{\linewidth}
\begin{lstlisting}[language=C, caption={transaction\_service.c --- A method
with business logic. The naming convention \texttt{transaction\_service\_*}
is the only thing binding this function to its type.}]
bool transaction_service_record_transaction(
    TransactionService             *svc,
    const RecordTransactionCommand *cmd,
    TransactionId *out_id, ServiceError *err)
{
    DomainError derr;
    Transaction *tx = transaction_create(..., &derr);
    if (!tx) { if (err) *err = service_error_from_domain(&derr);
                return false; }
    bool ok = svc->tx_repo.vtable
        ->insert(svc->tx_repo.ctx, tx, &new_id, &rerr);
        // -> dispatches through vtable
    transaction_free(tx);
    return ok;
}
\end{lstlisting}\end{minipage}

Rust formalises this relationship through \texttt{impl} blocks. All
functions declared inside an \texttt{impl} block belong to that type and
are only callable on an instance of it. Associated functions --- those
without \texttt{self} --- serve as constructors and are called with
\texttt{Type::function()} syntax:

\begin{lstlisting}[language=Rust, caption={transaction.rs (service) ---
\texttt{new} is an associated function: it belongs to the type's namespace
but does not require an existing instance. This is the idiomatic constructor
pattern in Rust.}]
impl<T, A> TransactionService<T, A>
where
    T: TransactionRepository,
    A: AccountRepository,
{
    pub fn new(tx_repo: Arc<T>, account_repo: Arc<A>) -> Self {
        Self { tx_repo, account_repo }
    }
}
\end{lstlisting}

Methods that operate on an existing instance use \texttt{self} as their
first parameter. The first parameter is not a convention --- it is a
language keyword with precise semantics: \texttt{\&self} grants read-only
access, \texttt{\&mut self} grants exclusive mutable access, and plain
\texttt{self} consumes the value entirely:

\noindent\begin{minipage}{\linewidth}
\begin{lstlisting}[language=Rust, caption={transaction.rs (service) ---
\texttt{\&self} is enforced by the borrow checker. There is no way to call
\texttt{record\_transaction} without a valid receiver, and no way to mutate
\texttt{self} through a read-only borrow.}]
// inside impl<T, A> TransactionService<T, A>
    pub async fn record_transaction(
        &self,          // read-only borrow -- compiler enforced
        cmd: RecordTransactionCommand,
    ) -> ServiceResult<TransactionId> {
        if cmd.value_cents <= 0 {
            return Err(ServiceError::Validation(
                "transaction value must be strictly positive".into()));
        }
        let tx = Transaction::create(
            cmd.from_account_id, cmd.to_account_id,
            cmd.transaction_type, cmd.value_cents,
            cmd.description,
        )?;
        let id = self.tx_repo.insert(&tx).await?;
        Ok(id)
    }
\end{lstlisting}\end{minipage}

Palmieri~\cite{Palmieri2022} notes that dot-notation method calls on
\texttt{\&self} allow the borrow checker to verify at compile time that no
two mutable references to the same data coexist during any method call.
This guarantee is impossible in C, where the same struct can be passed to
multiple functions simultaneously with no compiler objection.

\section{Polymorphism and Dynamic Dispatch}
\label{sec:ch4sec3}

Polymorphism allows a single interface to operate on objects of different concrete
types. C achieves this through function pointer tables (vtables), manually
constructed and stored alongside or within each object. Rust provides the
\texttt{trait} abstraction, which the compiler can dispatch either statically
(monomorphisation) or dynamically (trait objects).

In the Gentlix Bank service layer, repository abstractions must be swappable: the
same service logic should work against a PostgreSQL implementation in production
and an in-memory fake during testing. The C implementation defines a vtable struct
holding function pointers, and each repository struct stores a pointer to its
vtable alongside a \texttt{void *ctx} pointing to the actual implementation:

\noindent\begin{minipage}{\linewidth}
\begin{lstlisting}[language=C, caption={transaction\_service.h --- VTable pattern.
Each method is a function pointer; \texttt{ctx} holds implementation-specific
state.}]
typedef struct TransactionRepositoryVTable {
    bool (*get_by_id)(void *ctx, TransactionId id,
                      Transaction **out, RepoError *err);
    bool (*insert)(void *ctx, const Transaction *tx,
                   TransactionId *out_id, RepoError *err);
    bool (*list_for_account)(void *ctx, AccountId account_id,
                             int64_t limit, int64_t offset,
                             Transaction ***out_txs,
                             size_t *out_count,
                             RepoError *err);
} TransactionRepositoryVTable;

typedef struct TransactionRepository {
    const TransactionRepositoryVTable *vtable;
    void *ctx;
} TransactionRepository;
\end{lstlisting}\end{minipage}

To create a repository from a concrete implementation, adapter functions forward
calls from the vtable signature to the actual repo functions:

\enlargethispage{1\baselineskip}
\begin{lstlisting}[language=C, caption={transaction\_service.c --- Adapter
functions bridge the generic vtable to the concrete repo. The pattern matches the
``Class'' structure.}]
static bool tx_repo_insert_adapter(
    void *ctx, const Transaction *tx,
    TransactionId *out_id, RepoError *err) {
    return transaction_repo_insert(
        (TransactionRepo *)ctx, tx, out_id, err);
}

static const TransactionRepositoryVTable TX_REPO_VTABLE = {
    tx_repo_get_by_id_adapter,
    tx_repo_insert_adapter,
    tx_repo_list_for_account_adapter
};

TransactionRepository tx_repository_from_repo(
    TransactionRepo *repo) {
    TransactionRepository r;
    r.vtable = &TX_REPO_VTABLE;
    r.ctx    = repo;
    return r; }
\end{lstlisting}

This technique is described by Schreiner~\cite{Schreiner1993} as
\emph{dynamic linkage} --- a struct that pairs a vtable pointer with an
opaque \texttt{void *ctx} behaves exactly like a C++ object with a virtual
dispatch table, but constructed entirely by hand. The compiler cannot catch a 
vtable that is incorrectly populated or a function pointer with the wrong 
signature — both mistakes compile cleanly and only surface at runtime.

Rust's trait system formalises this pattern. A trait declares a set of
method signatures; any type that provides concrete definitions for all of
them automatically satisfies the interface. Unlike C's vtable, the compiler
verifies that every method signature matches exactly:

\begin{lstlisting}[language=Rust, caption={transaction.rs (service) ---
The trait defines the interface. \texttt{Send + Sync} bounds ensure that any
implementation can be shared safely across threads, replacing C's unverified
\texttt{void *ctx}.}]
#[async_trait]
pub trait TransactionRepository: Send + Sync {
    async fn get_by_id(&self, id: TransactionId)
        -> Result<Transaction, RepoError>;
    async fn insert(&self, tx: &Transaction)
        -> Result<TransactionId, RepoError>;
    async fn list_for_account(&self, account_id: AccountId,
        limit: i64, offset: i64)
        -> Result<Vec<Transaction>, RepoError>;
}

// Concrete implementation -- one impl block per type
#[async_trait]
impl TransactionRepository for TransactionRepo {
    async fn insert(&self, tx: &Transaction)
        -> Result<TransactionId, RepoError> {
        self.insert(tx).await   
        // delegates to the SQLx-backed method 
    }
    // get_by_id and list_for_account follow the same pattern
}
\end{lstlisting}

The same interface can then be used in two different ways. Static dispatch
--- triggered by a generic bound \texttt{T: TransactionRepository} --- lets
the compiler generate specialised code for each concrete type at compile
time, with no runtime overhead. Dynamic dispatch --- through
\texttt{Arc<dyn TransactionRepository>} --- creates a vtable at runtime,
identical in structure to C's manual approach but verified by the compiler:

\begin{lstlisting}[language=Rust, caption={Static dispatch vs dynamic
dispatch. The service uses generics for compile-time specialisation;
\texttt{dyn Trait} is available when the concrete type is unknown at
compile time.}]
// Static dispatch -- compiler generates one version per concrete T
pub struct TransactionService<T: TransactionRepository> {
    tx_repo: Arc<T>,   /* type known at compile time */
}

// Dynamic dispatch -- single version, type resolved at runtime
pub struct TransactionService {
    tx_repo: Arc<dyn TransactionRepository>,   // runtime vtable
}
\end{lstlisting}

Palmieri~\cite{Palmieri2022} recommends static dispatch when concrete types
are known at compile time, as it permits inlining and zero runtime cost.
In Gentlix Bank, the service layer uses generic bounds so that tests can
inject an in-memory fake repository while production uses the PostgreSQL
implementation --- the same flexibility as C's vtable swap, but checked
entirely at compile time.

\section{Inheritance and Composition}
\label{sec:ch4sec4}

Classical object-oriented languages use inheritance to reuse code by
extending existing types. Neither C nor Rust supports inheritance as a
built-in feature, but they each offer a different alternative.

In C, the standard approach is struct embedding. A ``derived'' struct places
a ``base'' struct as its first member, which makes pointer casting between
the two types safe and well-defined. Schreiner~\cite{Schreiner1993} illustrates
this with geometric shapes: a \texttt{Circle} embeds a \texttt{Point},
inheriting its coordinates without any language keyword. In Gentlix Bank,
the same idea appears in the service layer: \texttt{TransactionService}
embeds two repository structs directly by value, owning them as part of
its own memory layout:

\begin{lstlisting}[language=C, caption={transaction\_service.c ---
Composition through embedding. The service owns both repository structs
by value; no pointer indirection is needed to access their fields.}]
struct TransactionService {
    TransactionRepository tx_repo;      // embedded by value
    TxAccountRepository   account_repo; // embedded by value 
};
\end{lstlisting}

This provides full control over memory layout, but polymorphism is not included 
automatically — if dynamic dispatch is needed, vtables must be added on top, 
as described in Section ~\ref{sec:ch4sec3}.

Rust intentionally omits inheritance. The language designers argue that
inheritance couples child types tightly to parent implementations, making
refactoring harder and often breaking encapsulation. Instead, Rust promotes
composition of smaller value types and code sharing through traits.

In Gentlix Bank, domain entities follow this approach directly. A
\texttt{Transaction} is built from smaller validated value objects, each
enforcing its own invariants. The struct aggregates them without claiming
any ``is-a'' relationship:

\begin{lstlisting}[language=Rust, caption={transaction.rs (domain) ---
Composition over inheritance. Each field is a validated value type; the
compiler ensures that a \texttt{Transaction} can only exist with all
fields present and valid.}]
pub struct Transaction {
    id:               Option<TransactionId>,
    from_account_id:  AccountId,
    to_account_id:    AccountId,
    transaction_type: TransactionType,
    value_cents:      i64,
    description:      String,
    recorded_on:      DateTime<Utc>,
}
\end{lstlisting}

When behaviour must be shared across types, traits take the role that
abstract base classes play in other languages. A type can implement
multiple traits without the diamond problem that multiple inheritance
creates. Palmieri~\cite{Palmieri2022} notes that this separation --- types
hold data, traits define behaviour --- keeps each piece small and easy to
test independently, which is exactly the approach used across the Gentlix
Bank domain layer.

\section{Layered Architecture}
\label{sec:ch4sec5}

A layered architecture separates a backend application into distinct
responsibilities: domain logic, data access, application services, and HTTP
handling. Each layer communicates only with its immediate neighbour, so a
change in the database implementation does not affect the service layer, and
a change in the HTTP layer does not affect business rules.

\newpage
In C, these boundaries are enforced through header file separation and the
opaque pointer pattern from Section~\ref{sec:ch4sec1}. Each layer exposes
only an incomplete type and a set of functions; the layers below it are
invisible. The Gentlix Bank C backend follows this structure directly:

\begin{lstlisting}[language=C, caption={Header boundaries in the C backend.
Each layer sees only the interface of the layer below; concrete structs and
database handles never cross layer boundaries.}]
// domain/transaction.h -- business rules, no DB knowledge
typedef struct Transaction Transaction;
Transaction *transaction_create(AccountId from, AccountId to,
    TransactionType type, int64_t value_cents,
    const char *desc, DomainError *err);

// db/transaction_repo.h -- persistence, knows about PGconn
typedef struct TransactionRepo TransactionRepo;
TransactionRepo *transaction_repo_new(PGconn *db);

// service/transaction_service.h -- use-cases, no PGconn visible
typedef struct TransactionService TransactionService;
TransactionService *transaction_service_new(
    TransactionRepository tx_repo,
    TxAccountRepository   account_repo);
\end{lstlisting}

The service layer never sees a \texttt{PGconn *} or any SQLite handle.
It receives only a \texttt{TransactionRepository} --- a vtable and a
context pointer --- and calls methods through it. Concrete database
implementations are injected at startup by the application entry point,
which is the only place that knows about all layers simultaneously.

Rust enforces the same boundaries through module folders inside one Cargo
package. The \texttt{domain} module does not import \texttt{db}; the
\texttt{service} module defines repository traits; only \texttt{server}
wires HTTP handlers to concrete repositories at startup:

\begin{lstlisting}[language=Rust, caption={backend/rust/src/ module layout.
Folder boundaries replace C header boundaries; dependency direction is enforced
by \texttt{mod} declarations and imports.}]
backend/rust/src/
  domain/   -- entities and value objects
  db/       -- sqlx repositories
  service/  -- use-cases and repository traits
  server/   -- axum routes, JWT middleware, AppState
\end{lstlisting}

Dependency injection in the service layer happens through generic bounds.
The service does not import the \texttt{db} crate directly; it only requires
that its type parameter implements the repository trait:

\begin{lstlisting}[language=Rust, caption={service/transaction.rs ---
\texttt{TransactionService} depends on a trait, not a concrete type.}]
pub struct TransactionService<T, A>
where
    T: TransactionRepository, // any type implementing the trait
    A: AccountRepository,
{
    tx_repo:      Arc<T>,
    account_repo: Arc<A>,
}
\end{lstlisting}

Palmieri~\cite{Palmieri2022} describes this pattern as the key advantage
of trait-based dependency injection: the service is tested with an
in-memory fake, and the same code runs in production against PostgreSQL,
with no conditional compilation and no runtime cost from the swap.

\section*{Summary}

The five sections of this chapter show that C and Rust solve the same
design problems through fundamentally different mechanisms. In C, every
pattern is a convention: encapsulation through header boundaries,
methods through naming prefixes, polymorphism through hand-built vtables,
composition through struct embedding, and layered architecture through
include guards. These conventions work, but the compiler cannot verify
them. A misplaced pointer, a forgotten prefix, or an incorrectly populated
vtable all compile without error.

Rust replaces each convention with a language rule. Private fields,
\texttt{impl} blocks, traits, and module boundaries are not agreements
between developers --- they are constraints enforced at compile time. The
result is the same architecture, the same separation of concerns, and the
same flexibility for dependency injection, but with the compiler checking every 
access point automatically, something a human reviewer cannot reliably do.